% ── Chapter 1: Introduction ───────────────────────────────────
\chapter{บทนำ (Introduction)}

\section{ภูมิหลังและที่มาของปัญหา (Background and Problem Statement)}

การดื้อยาต้านจุลชีพ (Antimicrobial resistance หรือ AMR) ได้กลายเป็นหนึ่งในภัยคุกคามที่เร่งด่วนและซับซ้อนที่สุดต่อสาธารณสุขระดับโลกในศตวรรษที่ยี่สิบเอ็ด เมื่อแบคทีเรีย ไวรัส เชื้อรา หรือพยาธิ มีวิวัฒนาการกลไกในการทนทานต่อยาที่ออกแบบมาเพื่อกำจัดพวกมัน ทำให้การรักษามาตรฐานล้มเหลว และการติดเชื้อยังคงดำเนินต่อไป แพร่กระจาย และส่งผลให้มีการเสียชีวิตเพิ่มมากขึ้น องค์การอนามัยโลก (World Health Organization หรือ WHO) ได้จัดประเภท AMR เป็นลำดับความสำคัญระดับโลกที่สำคัญ โดยคาดการณ์ว่าการติดเชื้อที่ดื้อยาเป็นสาเหตุโดยตรงของการเสียชีวิตประมาณ 1.27 ล้านคนในปี 2019 เพียงปีเดียว และมีส่วนร่วมทางอ้อมต่อการเสียชีวิตเพิ่มเติมเกือบ 5 ล้านคนทั่วโลก \cite{who_amr2022} หากไม่มีการดำเนินการที่ประสานกัน มีการคาดการณ์ว่า AMR อาจคร่าชีวิตผู้คนได้ถึง 10 ล้านคนต่อปีภายในปี 2050 ซึ่งแซงหน้ามะเร็งในฐานะสาเหตุหลักของการเสียชีวิตทั่วโลก นอกเหนือจากความสูญเสียทางด้านมนุษยธรรมแล้ว AMR ยังก่อให้เกิดต้นทุนทางเศรษฐกิจที่มหาศาลผ่านการพักรักษาตัวในโรงพยาบาลที่ยาวนานขึ้น ยาในลำดับที่สองและสามที่มีราคาแพงกว่า และการลดลงของผลิตภาพของแรงงาน

หัวใจสำคัญในการต่อสู้กับ AMR คือการทดสอบความไวต่อยาต้านจุลชีพ (Antimicrobial Susceptibility Testing หรือ AST) ก่อนที่จะเริ่มการรักษาด้วยยาปฏิชีวนะ แพทย์ต้องกำหนดว่ายาชนิดใดที่ยังคงมีประสิทธิภาพต่อเชื้อก่อโรคแบคทีเรียเฉพาะที่ติดเชื้อในผู้ป่วย การสั่งจ่ายยาปฏิชีวนะที่สิ่งมีชีวิตนั้นดื้อยาอยู่แล้วไม่เพียงแต่ทำให้การรักษาล้มเหลว แต่ยังเป็นการคัดเลือกเชื้อที่ดื้อยามากขึ้นในประชากร ซึ่งทำให้วิกฤตการณ์ดำเนินต่อไป ดังนั้น การทำ AST ที่แม่นยำและทันเวลาจึงไม่ใช่เพียงแค่ขั้นตอนการวินิจฉัย แต่เป็นการแทรกแซงทางสาธารณสุข ความล่าช้าในการได้รับผลความไวต่อยา ซึ่งอาจใช้เวลาตั้งแต่ 24 ถึง 72 ชั่วโมงโดยใช้วิธีการเพาะเชื้อ บังคับให้แพทย์ต้องเริ่มการรักษาแบบครอบคลุมเชื้อกว้าง (empirical broad-spectrum therapy) ซึ่งเป็นการปฏิบัติที่เร่งการพัฒนาความดื้อยาด้วยตนเอง

วิธี Kirby-Bauer disk diffusion ซึ่งมีการอธิบายและตรวจสอบอย่างเป็นระบบครั้งแรกในทศวรรษ 1960 ยังคงเป็นเทคนิคที่ได้รับการยอมรับอย่างกว้างขวางที่สุดสำหรับ AST ในห้องปฏิบัติการจุลชีววิทยาทางคลินิกและห้องปฏิบัติการอ้างอิงทั่วโลก \cite{asm2009kirby,clsi2025} ความนิยมที่ยั่งยืนมาจากความเรียบง่าย ต้นทุนต่ำ และความน่าเชื่อถือเมื่อดำเนินการอย่างถูกต้อง ในวิธีนี้ จานอาหารเลี้ยงเชื้อ Mueller-Hinton agar จะถูกปลูกเชื้อด้วยสารแขวนลอยแบคทีเรียเป้าหมายมาตรฐาน (โดยทั่วไปจะปรับให้มีความขุ่นตามมาตรฐาน 0.5 McFarland) แผ่นดิสก์กระดาษที่ผลิตในเชิงพาณิชย์ ซึ่งแต่ละแผ่นจะชุ่มด้วยปริมาณยาปฏิชีวนะเฉพาะที่แม่นยำ จะถูกวางบนพื้นผิววุ้นที่ปลูกเชื้อไว้ที่ระยะห่างระหว่างแผ่นดิสก์ที่กำหนดเพื่อป้องกันการซ้อนทับของโซน หลังจากระยะเวลาการฟักตัวข้ามคืนที่ 35--37$^\circ$C ยาปฏิชีวนะจะแพร่ออกไปในแนวรัศมีจากแผ่นดิสก์แต่ละแผ่น สร้างระดับความเข้มข้นในวุ้น ในจุดที่ความเข้มข้นของยาปฏิชีวนะในพื้นที่สูงกว่าค่าความเข้มข้นต่ำสุดที่ยับยั้งเชื้อได้ (minimum inhibitory concentration หรือ MIC) ของแบคทีเรีย การเติบโตของแบคทีเรียจะถูกระงับ ส่งผลให้เห็นเป็นวงใสที่เรียกว่าโซนการยับยั้ง (zone of inhibition หรือ ZOI) ในจุดที่ความเข้มข้นลดลงต่ำกว่า MIC แบคทีเรียจะเติบโตตามปกติ

\begin{figure}[H]
  \centering
  \safeimg[width=0.85\textwidth]{figures/kirby_bauer_example.png}{Example of Kirby-Bauer disk diffusion plate showing zones of inhibition}
  \caption{Representative Kirby-Bauer disk diffusion plate illustrating zones of inhibition (ZOI)
    around antibiotic-impregnated disks. The diameter of each clearing zone is measured to
    determine antibiotic susceptibility.}
  \label{fig:kirby_bauer}
\end{figure}

เส้นผ่านศูนย์กลางของแต่ละ ZOI จะถูกวัดและเปรียบเทียบกับจุดตัดการตีความ (interpretive breakpoints) ที่ตีพิมพ์เพื่อจำแนกสิ่งมีชีวิตที่ทดสอบว่าเป็น ไวต่อยา (Susceptible: S), ระดับปานกลาง (Intermediate: I) หรือ ดื้อยา (Resistant: R) จุดตัดเหล่านี้ได้รับการดูแลและปรับปรุงอย่างเข้มงวดโดยองค์กรระหว่างประเทศหลักสองแห่ง ได้แก่ Clinical and Laboratory Standards Institute (CLSI) ในสหรัฐอเมริกา และ European Committee on Antimicrobial Susceptibility Testing (EUCAST) ในยุโรป \cite{eucast2023} CLSI ตีพิมพ์มาตรฐาน M02 ที่ควบคุมขั้นตอนการทดสอบ disk diffusion ในขณะที่ EUCAST ตีพิมพ์ตารางจุดตัดที่แก้ไขเป็นประจำทุกปี ซึ่งได้มาจากแบบจำลอง pharmacokinetic/pharmacodynamic (PK/PD), ค่าจุดตัดทางระบาดวิทยา (epidemiological cutoff values หรือ ECOFFs) และข้อมูลผลลัพธ์ทางคลินิก \cite{giske2022} ทั้งสององค์กรระบุเกณฑ์เส้นผ่านศูนย์กลางเฉพาะสำหรับสายพันธุ์และยาที่ต้องนำไปใช้อย่างถูกต้องเพื่อให้ได้ผลลัพธ์ที่มีความหมายทางคลินิก

แม้จะดูเรียบง่ายในการวัดเส้นผ่านศูนย์กลางของวงกลม แต่ในทางปฏิบัติ การวัด ZOI ด้วยตนเองนั้นมีแนวโน้มที่จะเกิดข้อผิดพลาดได้มาก ช่างเทคนิคในห้องปฏิบัติการใช้ไม้บรรทัด เวอร์เนียร์คาลิปเปอร์ หรืออุปกรณ์เทมเพลตเฉพาะในการวัดแต่ละโซน และความแปรปรวนระหว่างผู้สังเกต (inter-observer variability) เป็นความท้าทายที่ได้รับการยอมรับและยังคงมีอยู่ในสาขานี้ \cite{humphries2018} ปัจจัยต่างๆ เช่น โซนที่ทับซ้อนกันบางส่วนระหว่างแผ่นดิสก์ที่อยู่ติดกัน การเติบโตแบบไม่เป็นเนื้อเดียวกันหรือแบบ ``swarming'' ภายในขอบเขตโซน พื้นผิววุ้นที่ไม่สม่ำเสมอ การปลูกเชื้อบนจานที่ไม่สม่ำเสมอ การฟักตัวที่ไม่เหมาะสม และมุมการมองหรือการถ่ายภาพที่ไม่สอดคล้องกัน ล้วนส่งผลต่อความแปรปรวนในการวัด แม้แต่ความคลาดเคลื่อนเพียงมิลลิเมตรเดียวก็สามารถตัดสินผลทางคลินิกได้ ตัวอย่างเช่น CLSI 2025 กำหนดเส้นผ่านศูนย์กลาง ZOI ที่ 18\,mm หรือมากกว่าสำหรับ \textit{Staphylococcus epidermidis} ที่ทดสอบกับ oxacillin (1\,$\mu$g) ว่าเป็น ไวต่อยา ในขณะที่ 17\,mm หรือน้อยกว่าจะถูกจัดเป็น ดื้อยา เนื่องไม่มีหมวดหมู่ระดับปานกลางสำหรับคู่ยาและสิ่งมีชีวิตนี้ ทุกมิลลิเมตรจึงมีความหมาย นอกจากนี้ ห้องปฏิบัติการที่ประมวลผลกรณีจำนวนมากต้องการให้ช่างเทคนิควัดโซนหลายสิบโซนต่อกะการทำงาน ซึ่งทำให้ความเหนื่อยล้ากลายเป็นแหล่งที่มาของข้อผิดพลาดที่เพิ่มมากขึ้น

แพลตฟอร์ม AST อัตโนมัติที่มีอยู่ในปัจจุบัน เช่น ระบบ bioM\'{e}rieux VITEK 2, BD Phoenix และ Microscan WalkAway ช่วยแก้ปัญหาความท้าทายในการวัดเหล่านี้ผ่านฮาร์ดแวร์เฉพาะและซอฟต์แวร์ที่เป็นกรรมสิทธิ์ ระบบเหล่านี้ให้ผลลัพธ์ที่รวดเร็วและทำซ้ำได้ และมีการใช้กันอย่างกว้างขวางในห้องปฏิบัติการโรงพยาบาลขนาดใหญ่ อย่างไรก็ตาม ระบบเหล่านี้ต้องการการลงทุนจำนวนมาก---มักจะเป็นเงินหลายหมื่นดอลลาร์สหรัฐต่อยูนิต---รวมถึงค่าสารเคมีและค่าบำรุงรักษาอย่างต่อเนื่อง ซึ่งเป็นข้อจำกัดสำหรับคลินิกขนาดเล็ก โรงพยาบาลชุมชน และห้องปฏิบัติการในประเทศที่มีรายได้ต่ำและปานกลาง (low- and middle-income countries หรือ LMICs) สภาพแวดล้อมที่มีทรัพยากรจำกัดเหล่านี้มักเป็นที่ที่ได้รับภาระหนักที่สุดจาก AMR ทำให้เกิดความเหลื่อมล้ำอย่างชัดเจนระหว่างห้องปฏิบัติการที่ต้องการ AST ที่แม่นยำมากที่สุดกับห้องปฏิบัติการที่สามารถจ่ายค่าโซลูชันอัตโนมัติได้

ดังนั้น จึงมีเหตุผลที่น่าสนใจในการพัฒนาเครื่องมือซอฟต์แวร์ที่ขับเคลื่อนด้วย AI ราคาประหยัดที่สามารถวิเคราะห์รูปถ่ายมาตรฐานของจาน Kirby-Bauer AST งานวิจัยล่าสุดได้แสดงให้เห็นว่า machine learning และ artificial intelligence สามารถมีบทบาทสำคัญในการทำนายและต่อสู้กับการดื้อยาต้านจุลชีพ \cite{bilal2025} และความก้าวหน้าในการทดสอบความไวต่อยาที่ขับเคลื่อนด้วย AI กำลังขยายขอบเขตในทางปฏิบัติของ AST อัตโนมัติอย่างรวดเร็ว \cite{liao2025} เครื่องมือดังกล่าวสามารถติดตั้งบนฮาร์ดแวร์ทั่วไป เช่น แล็ปท็อปหรือแม้แต่สมาร์ทโฟน และจะไม่ต้องการเครื่องมือเฉพาะทางนอกเหนือจากจานวุ้นเอง ซึ่งอาจช่วยขยายขอบเขตของ AST ที่แม่นยำไปยังสภาพแวดล้อมที่ปัจจุบันต้องพึ่งพาการอ่านด้วยตนเอง รวมถึงสภาพแวดล้อมที่ได้รับคำแนะนำจากมาตรฐานห้องปฏิบัติการระดับชาติ เช่น มาตรฐานที่ดูแลโดยกรมวิทยาศาสตร์การแพทย์ สถาบันวิจัยวิทยาศาสตร์สาธารณสุข (NIH Thailand) \cite{dmsc_nih} โครงงานปัจจุบันนี้ตอบสนองความต้องการนี้โดยการพัฒนา ประเมิน และติดตั้งระบบที่ใช้ AI สำหรับการตรวจจับและวัด ZOI อัตโนมัติ

\section{แรงจูงใจ (Motivation)}

ความก้าวหน้าล่าสุดใน deep learning และ computer vision ได้สร้างโมเดลที่มีประสิทธิภาพเทียบเท่าหรือเหนือกว่ามนุษย์ในงานการจดจำรูปภาพและการตรวจจับวัตถุที่หลากหลาย สิ่งที่มีความเกี่ยวข้องเป็นพิเศษกับการวิเคราะห์ภาพ AST คือความก้าวหน้าในการตรวจจับและระบุตำแหน่งของโครงสร้างวงกลมหรือวงรีในพื้นหลังภาพที่ซับซ้อนและไม่เป็นระเบียบ แตกต่างจากชุดข้อมูลภาพธรรมชาติที่วัตถุมีความหลากหลายของรูปร่างอย่างมหาศาล จาน AST นำเสนอขอบเขตภาพที่มีข้อจำกัด: โซนวงกลมบนพื้นหลังวุ้นที่มีพื้นผิว พร้อมด้วยวัตถุอ้างอิงที่มีขนาดที่ทราบ (แผ่นดิสก์ยาปฏิชีวนะ) ที่ปรากฏในทุกภาพ โครงสร้างขอบเขตนี้ถูกนำมาใช้ประโยชน์โดยการเลือกการออกแบบของโครงงานปัจจุบัน

แรงจูงใจสำหรับโครงงานนี้เกิดขึ้นจากการสังเกตคู่ขนานสองประการ ประการแรก โมเดลการตรวจจับวัตถุที่ใช้ deep learning---โดยเฉพาะตัวตรวจจับแบบสองขั้นตอน เช่น Faster R-CNN และตัวตรวจจับแบบขั้นตอนเดียว เช่น YOLOv11---ได้แสดงความสามารถในการเรียนรู้ที่จะระบุโซนการยับยั้งจากข้อมูลการฝึกอบรมที่ติดป้ายกำกับ แม้ว่าโซนจะทับซ้อนกันบางส่วนหรือแสดงสัณฐานวิทยาที่ผิดปกติก็ตาม ประการที่สอง โซลูชันซอฟต์แวร์ที่ครอบคลุมสำหรับการอ่าน AST ไม่เพียงแต่ต้องตรวจจับโซนเท่านั้น แต่ยังต้องเชื่อมโยงแต่ละโซนเข้ากับแผ่นดิสก์ยาปฏิชีวนะที่เกี่ยวข้องอย่างถูกต้องด้วย หากไม่ทราบว่ายาปฏิชีวนะชนิดใดสร้างโซนใด การวัดเส้นผ่านศูนย์กลางก็จะไม่มีความหมายทางคลินิก ปัญหาการเชื่อมโยงนี้ต้องการการอ่านฉลากตัวอักษรและตัวเลขที่พิมพ์บนแผ่นดิสก์ยาปฏิชีวนะแต่ละแผ่น---ซึ่งเป็นงานที่อยู่ในขอบเขตของการจดจำอักขระด้วยแสง (optical character recognition หรือ OCR)

การรวมกันของกระดูกสันหลังการตรวจจับวัตถุสำหรับการระบุตำแหน่ง ZOI, เครื่องยนต์ OCR สำหรับการอ่านฉลากยาปฏิชีวนะ และฐานข้อมูลจุดตัดที่มีโครงสร้างสำหรับการจำแนกทางคลินิก ก่อให้เกิดรากฐานแนวคิดของระบบที่พัฒนาขึ้นที่นี่ องค์ประกอบแต่ละส่วนนำเสนอความท้าทายทางเทคนิคของตนเอง การตรวจจับวัตถุต้องรับมือกับสัญญาณรบกวนในภาพ แสงที่แปรปรวน และการทับซ้อนของโซน OCR ต้องอ่านข้อความขนาดเล็กที่มีความคมชัดต่ำบนพื้นผิวแผ่นดิสก์ที่มีความโค้ง ซึ่งอาจถูกถ่ายภาพในทิศทางใดก็ได้ ระบบย่อยการปรับเทียบ (calibration) ต้องแปลงการวัดมาตราส่วนพิกเซลเป็นมิลลิเมตรทางกายภาพโดยใช้ขนาดทางกายภาพที่ทราบของแผ่นดิสก์ยาปฏิชีวนะเป็นข้อมูลอ้างอิงภายใน

แรงจูงใจเพิ่มเติมคือโอกาสในการเปรียบเทียบโดยตรงและอย่างเข้มงวดระหว่างแนวทางการประมวลผลภาพแบบคลาสสิกและวิธีการ deep learning สมัยใหม่บนชุดข้อมูลและโปรโตคอลการประเมินเดียวกัน Hough Circle Transform ซึ่งเป็นเทคนิคที่จัดตั้งขึ้นเป็นอย่างดีสำหรับการตรวจจับโครงสร้างวงกลม ให้เกณฑ์มาตรฐานที่มีหลักการซึ่งโมเดล deep learning สามารถนำมาประเมินได้ การเปรียบเทียบดังกล่าวมอบหลักฐานเชิงประจักษ์เกี่ยวกับประโยชน์เชิงปฏิบัติของ deep learning สำหรับงานเฉพาะนี้ นอกเหนือจากที่มีให้จากเกณฑ์มาตรฐานเอนกประสงค์

\section{วัตถุประสงค์ (Objectives)}

โครงงานนี้มีวัตถุประสงค์หลักสี่ประการ:
\begin{enumerate}
  \item เพื่อศึกษาและดำเนินการตีความผลการทดสอบ Kirby-Bauer disk diffusion ตามมาตรฐาน CLSI 2025 \cite{clsi2025} และ EUCAST 2023 \cite{eucast2023} เพื่อให้มั่นใจว่าการจำแนกความไวต่อยาของระบบมีพื้นฐานทางคลินิกที่มั่นคง
  \item เพื่อพัฒนาและฝึกฝนโมเดล AI---โดยเฉพาะ Faster R-CNN และ YOLOv11n---ที่สามารถตรวจจับโซนการยับยั้งและวัดเส้นผ่านศูนย์กลางจากภาพถ่ายจาน AST ด้วยค่าความคลาดเคลื่อนสมบูรณ์เฉลี่ย (mean absolute error หรือ MAE) ที่หรือต่ำกว่าเกณฑ์ที่ยอมรับได้ทางคลินิกที่ 3\,mm \cite{humphries2018}
  \item เพื่อเปรียบเทียบความแม่นยำ ความเร็ว และความทนทานของแนวทางการตรวจจับสามแบบอย่างเข้มงวด: Faster R-CNN (deep learning แบบสองขั้นตอน), YOLOv11n (deep learning แบบขั้นตอนเดียว) และ Hough Circle Transform (เกณฑ์มาตรฐานการประมวลผลภาพแบบคลาสสิก) เพื่อเป็นฐานเชิงปริมาณในการเลือกโมเดลที่มีประสิทธิภาพดีที่สุดสำหรับการติดตั้งใช้งาน
  \item เพื่อติดตั้งใช้งานระบบเป็นเว็บแอปพลิเคชันที่ใช้งานได้จริงและหันหน้าเข้าหาผู้ใช้ ซึ่งเหมาะสมสำหรับการใช้งานจริงในสภาพแวดล้อมห้องปฏิบัติการทางคลินิก โดยผสานรวมโมเดลการตรวจจับเข้ากับไปป์ไลน์ OCR และฐานข้อมูลการจำแนกจุดตัดให้เป็นเครื่องมือที่สอดคล้องและเข้าถึงได้
\end{enumerate}

\section{ขอบเขตของโครงงาน (Project Scope)}

ขอบเขตของโครงงานนี้ถูกกำหนดและมีขอบเขตดังนี้:

\begin{itemize}
  \item \textbf{ข้อมูลนำเข้า (Input data):} ระบบยอมรับภาพถ่ายของจานวุ้น Mueller-Hinton ที่ใช้ในการทดลอง Kirby-Bauer AST มาตรฐาน คาดว่าภาพจะถูกถ่ายจากมุมมองจากบนลงล่าง พร้อมด้วยแสงที่สม่ำเสมอและระยะห่างระหว่างกล้องกับจานที่คงที่เพื่อรักษาความสม่ำเสมอของมาตราส่วน ภาพอาจมาจากกล้องในห้องปฏิบัติการ กล้องจุลทรรศน์ทางคลินิก หรือกล้องสมาร์ทโฟนมาตรฐาน

  \item \textbf{โมเดลการตรวจจับ (Detection models):} โครงงานนี้ฝึกฝน ประเมิน และเปรียบเทียบแนวทางการตรวจจับสามแบบ: Faster R-CNN ที่มีกระดูกสันหลังการสกัดคุณลักษณะ ResNet-50, YOLOv11n (รุ่น nano ของ Ultralytics) และ Hough Circle Transform ที่ดำเนินการผ่าน OpenCV แต่ละโมเดลจะได้รับการประเมินบนชุดข้อมูลทดสอบที่แยกไว้ชุดเดียวกันภายใต้เงื่อนไขเดียวกัน

  \item \textbf{เกณฑ์การประเมิน (Evaluation criteria):} ตัวชี้วัดความแม่นยำหลักคือค่าความคลาดเคลื่อนสมบูรณ์เฉลี่ย (MAE) ในหน่วยมิลลิเมตรระหว่างเส้นผ่านศูนย์กลาง ZOI ที่ระบบทำนายและเส้นผ่านศูนย์กลางอ้างอิง (ground-truth) ที่วัดโดยนักจุลชีววิทยาที่ผ่านการฝึกอบรม ตัวชี้วัดรอง ได้แก่ ความเร็วในการประมวลผล (เฟรมต่อวินาที) และการระลึกในการตรวจจับ (detection recall) เป้าหมาย MAE ที่ยอมรับได้ทางคลินิกคือต่ำกว่า 3\,mm

  \item \textbf{การปรับเทียบ (Calibration):} การแปลงพิกเซลเป็นมิลลิเมตรดำเนินการโดยใช้แผ่นดิสก์ยาปฏิชีวนะเป็นตัวอ้างอิงภายในภาพ เนื่องจากเส้นผ่านศูนย์กลางทางกายภาพของแผ่นดิสก์ยาปฏิชีวนะมาตรฐานถูกกำหนดอย่างแม่นยำไว้ที่ 6\,mm ตามข้อกำหนดของ CLSI ไม่จำเป็นต้องมีเป้าหมายการปรับเทียบภายนอกหรือไม้บรรทัดในมุมมองของภาพ

  \item \textbf{เว็บแอปพลิเคชัน (Web application):} เว็บแอปพลิเคชันขั้นสุดท้ายช่วยให้ผู้ใช้สามารถอัปโหลดภาพจากไฟล์หรือถ่ายภาพแบบเรียลไทม์ผ่านกล้องที่เชื่อมต่อ ระบบจะส่งคืนแดชบอร์ดที่แสดงโซนที่ตรวจพบ ฉลากยาปฏิชีวนะที่ได้รับมอบหมาย (จาก OCR) เส้นผ่านศูนย์กลางที่วัดได้ และการจำแนก S/I/R ขั้นสุดท้ายสำหรับคู่ยาปฏิชีวนะและแผ่นดิสก์แต่ละคู่ตามมาตรฐาน CLSI หรือ EUCAST ที่เลือก ผลลัพธ์ย้อนหลังจะถูกจัดเก็บไว้สำหรับการตรวจสอบและทบทวน

  \item \textbf{นอกเหนือขอบเขต (Out of scope):} ระบบไม่ได้ดำเนินการระบุสายพันธุ์หรือชนิดของแบคทีเรีย ผู้ใช้ต้องระบุสิ่งมีชีวิตที่กำลังทดสอบ (เช่น \textit{Staphylococcus aureus}, \textit{Escherichia coli}) เพื่อให้จุดตัดที่ถูกต้องถูกนำไปใช้ ระบบไม่รองรับการทดสอบ broth microdilution, E-test หรือรูปแบบ AST อื่นๆ ผลลัพธ์จากเครื่องมือซอฟต์แวร์นี้มีวัตถุประสงค์เพื่อช่วยเหลือ ไม่ใช่เพื่อแทนที่การตัดสินใจของนักจุลชีววิทยาที่มีคุณสมบัติเหมาะสม
\end{itemize}

\section{ผลที่คาดว่าจะได้รับ (Expected Outcomes)}

ผลลัพธ์ที่เป็นรูปธรรมที่คาดหวังของโครงงานนี้มีดังนี้:

\begin{enumerate}
  \item \textbf{ไปป์ไลน์การตรวจจับด้วย AI (AI detection pipeline):} ไปป์ไลน์อัตโนมัติแบบครบวงจรที่สามารถตรวจจับและวัดเส้นผ่านศูนย์กลาง ZOI จากภาพจาน AST โมเดลที่มีประสิทธิภาพดีที่สุดคาดว่าจะบรรลุค่า MAE ต่ำกว่า 3\,mm ในชุดข้อมูลทดสอบที่แยกไว้ ซึ่งสอดคล้องกับเกณฑ์การยอมรับทางคลินิกที่กำหนดโดย Humphries et al.\ \cite{humphries2018}

  \item \textbf{การวิเคราะห์โมเดลเชิงเปรียบเทียบ (Comparative model analysis):} การเปรียบเทียบเชิงปริมาณที่เข้มงวดของ Faster R-CNN, YOLOv11n และ Hough Circle Transform บนชุดข้อมูลเดียวกัน โดยบันทึกข้อดีข้อเสียระหว่างความแม่นยำในการตรวจจับ ความเร็วในการประมวลผล และความซับซ้อนในการดำเนินการสำหรับขอบเขต AST การวิเคราะห์นี้มีวัตถุประสงค์เพื่อให้ข้อมูลแก่ผู้ปฏิบัติงานในอนาคตในการเลือกแนวทางการตรวจจับสำหรับงานวิเคราะห์ภาพทางการแพทย์ที่คล้ายคลึงกัน

  \item \textbf{เว็บแอปพลิเคชันที่ติดตั้งใช้งาน (Deployed web application):} เว็บแอปพลิเคชันที่ทำงานได้อย่างสมบูรณ์ซึ่งรวมโมเดลการตรวจจับที่มีประสิทธิภาพดีที่สุด ไปป์ไลน์การจดจำฉลากยาปฏิชีวนะที่ใช้ EasyOCR และฐานข้อมูลการจำแนกจุดตัด CLSI 2025/EUCAST 2023 มอบเวิร์กโฟลว์แบบครบวงจรตั้งแต่การอัปโหลดภาพไปจนถึงรายงานความไวต่อยา

  \item \textbf{ชุดข้อมูลที่มีคำอธิบายประกอบ (Annotated dataset):} ชุดข้อมูลที่คัดสรรและติดป้ายกำกับของภาพจาน AST จริงจำนวน 287 ภาพ ซึ่งขยายเป็น 543 ภาพผ่านการเพิ่มพูนข้อมูล (augmentation) และเสริมด้วยภาพสังเคราะห์ 2,000 ภาพ เป็นทรัพยากรที่สามารถนำกลับมาใช้ใหม่ได้สำหรับการวิจัยในอนาคตเกี่ยวกับ AST ที่ได้รับความช่วยเหลือจาก AI
\end{enumerate}

\section{โครงสร้างของรายงาน (Report Structure)}

ส่วนที่เหลือของรายงานนี้จัดลำดับดังนี้ บทที่ 2 ให้ทฤษฎีพื้นฐานและการทบทวนงานวิจัยที่เกี่ยวข้อง ครอบคลุมวิธี Kirby-Bauer AST และมาตรฐานทางคลินิก สถาปัตยกรรมและหลักการของโมเดลการตรวจจับวัตถุที่ใช้ (Faster R-CNN, YOLOv11n และ Hough Circle Transform), เครื่องยนต์ OCR ที่ใช้สำหรับการจดจำฉลากยาปฏิชีวนะ, เทคนิคการเตรียมการประมวลผลภาพที่สำคัญ และการทบทวนงานที่ตีพิมพ์ก่อนหน้านี้เกี่ยวกับการวิเคราะห์ภาพ AST อัตโนมัติ บทที่ 3 รายละเอียดระเบียบวิธีวิจัย รวมถึงเวิร์กโฟลว์การรวบรวมและเตรียมชุดข้อมูลทั้งหมด ขั้นตอนการติดป้ายกำกับ โปรโตคอลการฝึกโมเดล แนวทางการปรับเทียบ การออกแบบไปป์ไลน์ OCR ตรรกะการจำแนกความไวต่อยา และสถาปัตยกรรมระบบโดยรวม บทที่ 4 นำเสนอผลการทดลอง รวมถึงตัวชี้วัดความแม่นยำในการตรวจจับ การทดสอบความเร็ว ประสิทธิภาพของ OCR และการวิเคราะห์แหล่งที่มาของข้อผิดพลาดและการปรับแต่งการปรับเทียบ บทที่ 5 สรุปรายงานด้วยการสรุปสิ่งที่ค้นพบ ข้อจำกัด และแนวทางสำหรับงานในอนาคต
